\documentclass{article}%
\usepackage[T1]{fontenc}%
\usepackage[utf8]{inputenc}%
\usepackage{lmodern}%
\usepackage{textcomp}%
\usepackage{lastpage}%
\usepackage{geometry}%
\geometry{margin=2.2cm}%
%
\usepackage{amsmath}%
\usepackage{amssymb}%
\usepackage{booktabs}%
\usepackage{longtable}%
\usepackage{array}%
\usepackage{hyperref}%
\hypersetup{hidelinks}%
\title{ITA 2027 {-} Física (sem judge) x Poliedro\\[4pt]\large Comparação com o gabarito Poliedro Resolve (ITA 2027, 2ª fase, dia 3 {-} Física)\thanks{Gerado com gpt-resolve: \url{https://github.com/lgabs/gpt-resolve}}}%
\date{30/10/2026}%
%
\begin{document}%
\normalsize%
\maketitle%
\section*{Resumo}\textbf{Nota média: 10,0/10} (24 de 24 itens com a resolução correta; cada item vale o mesmo dentro da questão).\par\medskip%
\begin{longtable}{@{}c c p{0.78\textwidth}@{}}\toprule Questão & Nota & Resumo \\ \midrule
1 & 10,0 & A resolução está correta nos três itens e coincide com o gabarito do Poliedro. (a) k\_c = MgL/2. (b) Para k/k\_c $\geq$ 1, só $\theta$ = 0, estável. Para 2/$\pi$ < k/k\_c < 1, três equilíbrios: $\theta$ = 0 instável e $\pm$$\theta$\_e estáveis. Para k/k\_c $\leq$ 2/$\pi$, só $\theta$ = 0, instável. (c) U = k$\theta$$^2$/2 + (MgL/2)cos$\theta$, com curvas em U, em W e em ? conforme o caso. Refiz as contas numericamente (raízes de r$\theta$ = sen$\theta$, $\theta$\_e = 1,2757 para r = 0,75, U nos extremos) e tudo confere. A resolução ainda traz justificativas a mais que o gabarito, todas corretas: o mínimo quártico em k = k\_c, U'' nos pontos $\pm$$\theta$\_e e os valores de U nos extremos. \\
2 & 10,0 & A resolução da questão 2 está correta e confere com o gabarito do Poliedro nos dois itens.\newline%
\newline%
{-} **Item a:** a resolução dá H\_n = 0,49\^{}n$\cdot$V$^2$/(8g). É a mesma expressão do gabarito, porque 0,49\^{}n = 0,7\^{}(2n).\newline%
{-} **Item b:** os dois chegam a $\Delta$t = 10V/(3g), somando a PG dos tempos de voo (de razão e = 0,7).\newline%
\newline%
O enunciado transcrito confere com o PDF original. Refiz as contas em Python e os resultados batem. Não há divergência. Ficam só observações de pouco peso:\newline%
{-} a escolha do instante inicial do intervalo no item b;\newline%
{-} a idealização de infinitos choques em tempo finito;\newline%
{-} erros de digitação na notação da resolução do gabarito, que não alteram o resultado. \\
3 & 10,0 & Nos dois itens, as respostas da resolução e do gabarito são algebricamente equivalentes e estão corretas. No item a, D = R$\surd$(1 + (2h/g)$\surd$($\mu$e$^2$g$^2$/R$^2$ - $\alpha$$^2$)), escrito de outra forma. No item b, t\_p = $\omega$\_e/$\alpha$ + $\surd$(2h/g) + $\omega$\_e$\cdot$R/($\mu$c g), com $\omega$\_e = ($\mu$e$^2$g$^2$/R$^2$ - $\alpha$$^2$)\^{}\{1/4\}. Conferi numericamente e com uma simulação direta: dá t1 = 2,753 s, D = 2,420 m e t\_p = 4,185 s, iguais aos valores de ordem de grandeza da resolução. A resolução declara bem as hipóteses ($\mu$e$\cdot$g > $\alpha$R e velocidade horizontal mantida no choque). Já a resolução detalhada do gabarito tem erros de digitação, como $\mu$e$\cdot$g sem o quadrado dentro das raízes, mas eles não chegam às respostas resumidas por item. \\
4 & 10,0 & A resolução da questão 4 está correta e bate com o gabarito do Poliedro. No item a, as duas dão h\_f = h\_i$\cdot$((m+M)/m)\^{}\{3/5\}. No item b, a equação da resolução é algebricamente idêntica à interpretação 1 do Poliedro (adiabática reversível): 3(m+M)h\_i{[}1 - (h\_i/h\_max)\^{}\{2/3\}{]} = 2m(h\_max - h\_i). Confirmei a equivalência simbólica e numericamente, assim como o exemplo M = m (h\_max $\approx$ 2,26$\cdot$h\_i) e a análise de unicidade da raiz. O Poliedro observa que o enunciado não basta para determinar p\_f; a resolução contorna isso declarando a hipótese de gás isentrópico, que é a leitura mais natural para 'altura máxima'. \\
5 & 10,0 & Na questão 5, a resolução bate com o gabarito do Poliedro nos dois itens, e os dois estão certos. No item a, p', a natureza virtual da imagem e as alturas coincidem; conferi pela equação do dioptro com R\_alg = -R e por um teste numérico. No item b, a resolução dá $\lambda$(D + n2p/n1)/(2n2h), que com p ? D é o $\lambda$D/(2n2h) do gabarito. Um traçado exato de caminho óptico confirmou as duas formas. O único ponto que vale levar a um professor é o enunciado não dizer se $\lambda$ é no vácuo ou no meio n1; as duas respostas assumem o vácuo. \\
6 & 10,0 & Os três itens da resolução batem com o gabarito do Poliedro. a) A carga transferida é q = 3C$\varepsilon$/4. b) A mola fica comprimida de 2d/3 e as placas terminam a d/3 uma da outra. c) A constante é k = 27C$\varepsilon$$^2$/(64d$^2$). Refiz as contas com frações exatas e elas conferem. A transcrição bate com o enunciado original (página 7 do PDF). Não há divergência de resultado. Os pontos levantados são todos de baixa importância: o equilíbrio é marginalmente estável, o efeito de pêndulo é desprezado nas duas resoluções e o texto extraído do gabarito está com a formatação quebrada. \\
7 & 10,0 & A resolução está correta e coincide com o gabarito do Poliedro nos dois itens. No item a, |V\_o| = (V\_s/4)|G$\varepsilon$+$\alpha$$\Delta$T|. No item b, a ponte completa com quatro extensômetros (R2 e R3 na tração, R1 e R4 na compressão, ou o inverso) dá V\_o = V\_sG$\varepsilon$ sem efeito térmico. Conferi a fórmula linear da ponte, o caso de um extensômetro em cada braço e a expressão exata V\_sG$\varepsilon$/(1+$\alpha$$\Delta$T); todos batem. Não há divergência a apontar, só observações de baixa importância sobre a meia ponte como resposta parcial e a convenção de sinal de $\varepsilon$. \\
8 & 10,0 & A resolução está correta e bate com o gabarito do Poliedro nos três itens. Item a: v = $\surd$(2qV/m). Item b: movimento helicoidal uniforme, com R = m v sen$\theta$/(qB) e T = 2$\pi$m/(qB). Item c: cos$\theta$ = qBL/(2$\pi$nmv) = (BL/(2$\pi$n))$\surd$(q/(2mV)), com n inteiro positivo. A única diferença é a restrição cos$\theta$ < 1 na resolução contra cos$\theta$ $\leq$ 1 no gabarito, que é de redação e não altera a resposta. Conferi o enunciado no PDF original e refiz a verificação numérica: as duas formas da fórmula devolvem o cos$\theta$ usado no exemplo. \\
9 & 10,0 & Os três itens da resolução coincidem com o gabarito do Poliedro. No item a, E = E0{[}1+(-1)\^{}n{]} cos(kx-$\omega$t) $\hat z$, propagando{-}se no sentido +x. No item b, n par dá interferência construtiva (amplitude 2E0) e n ímpar dá interferência destrutiva total. No item c, B = -B0 cos(kx-$\omega$t) $\hat y$, com E$\times$B no sentido +x. Refiz a lei de Faraday e o sinal de B\_y está correto. O enunciado transcrito confere com o PDF original. Restam só observações menores: a leitura alternativa de B0 no item c e o caso trivial de n ímpar, em que o campo é nulo. \\
10 & 10,0 & A resolução está correta nos dois itens e concorda com o gabarito do Poliedro. No item a, as duas fórmulas de Compton são idênticas. No item b, a resolução dá |V| = c($\nu$$^2$ - $\nu$'$^2$)/($\nu$$^2$ + $\nu$'$^2$) = c$\cdot$x(2+x)/(x$^2$+2x+2), que é a mesma expressão c{[}(1+x)$^2$ - 1{]}/{[}(1+x)$^2$ + 1{]} do gabarito, com x = h$\nu$(1-cos$\theta$)/(m\_ec$^2$). As duas indicam o mesmo sentido: O vai ao encontro do fóton espalhado. Conferi numericamente a equivalência e a condição Doppler. A única falha está no gabarito: erros de digitação nas linhas intermediárias da dedução (quadrados e parênteses faltando), que não mudam a resposta final. \\
\bottomrule\end{longtable}%
\section*{Questão 1}%
\href{https://poliedroresolve.sistemapoliedro.com.br/vestibulares/ita/2027/ita-2-fase-dia-3-22-10-2026-2-fase-dia-3/dissertativa-1-fisica-2-fase-dia-3-ita-2027}{Resolução do gabarito}\par\smallskip%
\begin{longtable}{@{}p{0.06\textwidth} p{0.33\textwidth} p{0.33\textwidth} p{0.1\textwidth} p{0.08\textwidth}@{}}\toprule Item & Resolução & Gabarito & Veredito & Certo \\ \midrule
a & k\_c = MgL/2 & k\_c = MgL/2 & igual & ambos \\
b & k/k\_c $\geq$ 1: um equilíbrio ($\theta$=0, estável); 2/$\pi$ < k/k\_c < 1: três equilíbrios ($\theta$=0 instável; $\pm$$\theta$\_e estáveis, com sen$\theta$\_e/$\theta$\_e = k/k\_c); 0 < k/k\_c $\leq$ 2/$\pi$: um equilíbrio ($\theta$=0, instável) & k/k\_c $\geq$ 1: um único equilíbrio, $\theta$=0, estável; 2/$\pi$ < k/k\_c < 1: três equilíbrios, $\theta$=0 instável e $\pm$$\theta$0 estáveis (k$\theta$0 = k\_c sen$\theta$0); k/k\_c $\leq$ 2/$\pi$: um único equilíbrio ($\theta$=0), instável & igual & ambos \\
c & U\_total = k$\theta$$^2$/2 + (MgL/2)cos$\theta$; k $\geq$ k\_c: poço único em U com mínimo MgL/2 em 0; 2/$\pi$ < k/k\_c < 1: poço duplo em W com máximo local MgL/2 em 0 e mínimos em $\pm$$\theta$\_e; 0 < k/k\_c $\leq$ 2/$\pi$: curva em ? com máximo em 0, decrescendo até $\pi$$^2$k/8 em $\pm$$\pi$/2 & U($\theta$) = (MgL/2)cos$\theta$ + k$\theta$$^2$/2; k $\geq$ k\_c: poço único com mínimo em 0; 2/$\pi$ < k/k\_c < 1: duplo poço (máximo local em 0, mínimos em $\pm$$\theta$0); k/k\_c $\leq$ 2/$\pi$: máximo em 0 e U decrescendo para as bordas & igual & ambos \\
\bottomrule\end{longtable}%
\begin{itemize}%
\item%
\textbf{a)} O torque resultante é $\tau$ = -k$\theta$ + (MgL/2)sen$\theta$. Para $\theta$ pequeno, $\tau$ $\approx$ -(k - MgL/2)$\theta$, que só é restaurador se k > MgL/2. A resolução vai além do gabarito em dois pontos, ambos corretos: calcula $\omega$$^2$ = (k - MgL/2)/(ML$^2$/3) e mostra que em k = k\_c o mínimo é quártico (U - U(0) $\approx$ k\_c$\theta$$^{4}$/24), portanto ainda estável. Isso bate com o enunciado, segundo o qual a estabilidade se perde para k < k\_c.%
\item%
\textbf{b)} As faixas, as quantidades e as naturezas dos equilíbrios coincidem. A condição sen$\theta$\_e/$\theta$\_e = k/k\_c é a mesma que k$\theta$0 = k\_c sen$\theta$0. Conferi numericamente: para r = 1,5, 1, 0,4 e 2/$\pi$ a única raiz de r$\theta$ = sen$\theta$ em (-$\pi$/2, $\pi$/2) é $\theta$ = 0. Para r = 0,75 as raízes são 0 e $\pm$1,2757. A verificação de estabilidade U''($\theta$\_e) = k\_c cos$\theta$\_e (tan$\theta$\_e - $\theta$\_e)/$\theta$\_e > 0 está correta. A aproximação $\theta$\_e $\approx$ $\surd$(6(1-r)) também está correta.%
\item%
\textbf{c)} A expressão de U é a mesma e as três formas qualitativas também. Os valores extras da resolução estão corretos: U(0) = MgL/2; U($\pm$$\pi$/2) = $\pi$$^2$k/8 $\leq$ ($\pi$/4)(MgL/2) quando r $\leq$ 2/$\pi$; para r = 0,75, U($\theta$\_e)/k\_c = 0,9011 em $\theta$\_e = 1,2757, que são os valores usados no gráfico. A monotonicidade no caso r $\leq$ 2/$\pi$ vem de U' = k\_c$\theta$(r - sen$\theta$/$\theta$) < 0 para $\theta$ > 0. Não foi possível comparar ponto a ponto com o esboço do gabarito, que está em imagem não disponível.%
\end{itemize}%
\paragraph{Pontos de dúvida}%
\begin{itemize}%
\item%
\textit{(key)} O esboço do gabarito no item c está em uma imagem (figura2\_vetorial) que não está acessível aqui. A comparação foi feita só pela descrição textual (poço único, duplo poço, máximo em 0 decrescendo para as bordas).%
\item%
\textit{(statement)} Nos casos{-}limite k = k\_c e k/k\_c = 2/$\pi$ a classificação depende de convenções. Em k = k\_c, $\theta$ = 0 continua estável por ser um mínimo quártico, e o enunciado ('deixa de ser estável para k < k\_c') confirma isso. Em k/k\_c = 2/$\pi$ os pontos $\pm$$\theta$\_e caem em $\pm$$\pi$/2, fora do intervalo aberto. A resolução e o gabarito tratam os dois limites da mesma forma.%
\item%
\textit{(solution)} A resolução diz que para k $\leq$ MgL/$\pi$ a haste 'tomba sobre o suporte'. É uma interpretação física extra, não pedida, e não afeta a resposta. A energia elástica é tomada como zero em $\theta$ = 0, uma escolha implícita adotada também pelo gabarito.%
\end{itemize}%
\section*{Questão 2}%
\href{https://poliedroresolve.sistemapoliedro.com.br/vestibulares/ita/2027/ita-2-fase-dia-3-22-10-2026-2-fase-dia-3/dissertativa-2-fisica-2-fase-dia-3-ita-2027}{Resolução do gabarito}\par\smallskip%
\begin{longtable}{@{}p{0.06\textwidth} p{0.33\textwidth} p{0.33\textwidth} p{0.1\textwidth} p{0.08\textwidth}@{}}\toprule Item & Resolução & Gabarito & Veredito & Certo \\ \midrule
a & H\_n = (0,49)\^{}n V$^2$/(8g) & h\_n = (e\^{}n V sen30$^\circ$)$^2$/(2g) = 0,7\^{}(2n) V$^2$/(8g) & equivalente & ambos \\
b & $\Delta$t = V/{[}g(1-e){]} = 10V/(3g), contado desde o lançamento, com E\_min = 3mV$^2$/8 (só resta a energia cinética horizontal) & T = (2V sen30$^\circ$/g)$\cdot$1/(1-e) = 10V/(3g), após infinitos choques (velocidade vertical nula, sobra só a horizontal) & igual & ambos \\
\bottomrule\end{longtable}%
\begin{itemize}%
\item%
\textbf{a)} Mesma expressão, porque 0,7\^{}(2n) = 0,49\^{}n. Refazendo: a componente vertical logo após o n{-}ésimo choque é v\_y,n = e\^{}n$\cdot$V$\cdot$sen30$^\circ$ = e\^{}n$\cdot$V/2, e então H = v\_y,n$^2$/(2g) = e\^{}(2n)$\cdot$V$^2$/(8g). Conferi numericamente para n = 0\ldots{}4. As duas usam a mesma hipótese: o choque com o solo liso muda só a componente normal (vertical).%
\item%
\textbf{b)} A física e o resultado são os mesmos. Os tempos de voo t\_k = 2e\^{}k$\cdot$V$\cdot$sen30$^\circ$/g = (V/g)e\^{}k formam uma PG de razão 0,7. A soma é (V/g)/(1-0,7) = 10V/(3g); conferi numericamente (3,333\ldots{}$\cdot$V/g). A resolução também calcula corretamente a energia mínima, E\_min = (1/2)m($\surd$3V/2)$^2$ = 3mV$^2$/8, e a distância horizontal x = ($\surd$3/2)V$\cdot$10V/(3g) = 5$\surd$3V$^2$/(3g). O gabarito não pede esses dois valores, e eles estão certos.%
\end{itemize}%
\paragraph{Pontos de dúvida}%
\begin{itemize}%
\item%
\textit{(statement)} No item b, "intervalo de tempo" é lido como o tempo contado desde o lançamento até a energia mínima. Resolução e gabarito usam a mesma leitura, mas o enunciado não diz explicitamente qual é o instante inicial.%
\item%
\textit{(solution)} O mínimo de energia só é atingido no limite de infinitos choques, que somam um tempo finito (um paradoxo do tipo Zenão). Resolução e gabarito tratam esse limite como o instante em que o mínimo é atingido. Isso é o esperado, mas vale comentar que a validade do modelo de restituição constante para choques arbitrariamente fracos é uma idealização.%
\item%
\textit{(key)} Na resolução do gabarito há erros de digitação de notação: "(e\^{}nVsen$\theta$\^{}2)/2g", que deveria ser (e\^{}n V sen$\theta$)$^2$/(2g), e um "2gh" a mais na linha final do item a. Esses erros não mudam o resultado.%
\end{itemize}%
\section*{Questão 3}%
\href{https://poliedroresolve.sistemapoliedro.com.br/vestibulares/ita/2027/ita-2-fase-dia-3-22-10-2026-2-fase-dia-3/dissertativa-3-fisica-2-fase-dia-3-ita-2027}{Resolução do gabarito}\par\smallskip%
\begin{longtable}{@{}p{0.06\textwidth} p{0.33\textwidth} p{0.33\textwidth} p{0.1\textwidth} p{0.08\textwidth}@{}}\toprule Item & Resolução & Gabarito & Veredito & Certo \\ \midrule
a & D = $\surd$(R$^2$ + (2hR/g)$\cdot$$\surd$($\mu$e$^2$g$^2$ - $\alpha$$^2$R$^2$)), supondo $\mu$e$\cdot$g > $\alpha$R & D = R$\cdot$$\surd$(1 + (2h/g)$\cdot$$\surd$($\mu$e$^2$g$^2$/R$^2$ - $\alpha$$^2$)), com $\omega$\_e$^2$ = $\surd$($\mu$e$^2$g$^2$/R$^2$ - $\alpha$$^2$) & equivalente & ambos \\
b & t\_p = {[}($\mu$e g)$^2$ - ($\alpha$R)$^2${]}\^{}\{1/4\}/($\alpha$$\surd$R) + $\surd$(2h/g) + $\surd$R$\cdot${[}($\mu$e g)$^2$ - ($\alpha$R)$^2${]}\^{}\{1/4\}/($\mu$c g) & t\_parada = $\omega$\_e/$\alpha$ + $\surd$(2h/g) + $\omega$\_e$\cdot$R/($\mu$c g), com $\omega$\_e = ($\mu$e$^2$g$^2$/R$^2$ - $\alpha$$^2$)\^{}\{1/4\} & equivalente & ambos \\
\bottomrule\end{longtable}%
\begin{itemize}%
\item%
\textbf{a)} Basta pôr R$^2$ em evidência: (2hR/g)$\cdot$$\surd$($\mu$e$^2$g$^2$ - $\alpha$$^2$R$^2$) = (2hR$^2$/g)$\cdot$$\surd$($\mu$e$^2$g$^2$/R$^2$ - $\alpha$$^2$), e daí D = R$\cdot$$\surd$(1 + (2h/g)$\cdot$$\surd$($\mu$e$^2$g$^2$/R$^2$ - $\alpha$$^2$)), que é a fórmula do gabarito. A física é a mesma nas duas: o atrito estático chega a $\mu$e$\cdot$m$\cdot$g quando R$\surd$($\alpha$$^2$ + $\omega$$^{4}$) = $\mu$e$\cdot$g, a massa sai pela tangente com v = $\omega$\_e$\cdot$R, a queda dura $\surd$(2h/g) e Pitágoras dá a distância (o raio é perpendicular à tangente). Conferi numericamente com valores aleatórios (diferença \textasciitilde{}1e{-}15) e com uma simulação direta do disco: para R = 2 m, $\alpha$ = 0,5 rad/s$^2$, $\mu$e = 0,4 e h = 1,2 m, saem t1 = 2,753 s e D = 2,420 m, iguais aos da resolução.%
\item%
\textbf{b)} Temos $\omega$\_e = {[}($\mu$e g)$^2$ - ($\alpha$R)$^2${]}\^{}\{1/4\}/$\surd$R. Logo $\omega$\_e/$\alpha$ é igual ao t1 da resolução, e $\omega$\_e$\cdot$R/($\mu$c g) = $\surd$R$\cdot${[}($\mu$e g)$^2$ - ($\alpha$R)$^2${]}\^{}\{1/4\}/($\mu$c g) é igual ao t\_d. Os dois somam as mesmas três parcelas (instante do escorregamento, queda e frenagem com desaceleração $\mu$c$\cdot$g) e mantêm a velocidade horizontal no choque com o solo. A checagem numérica bate (diferença \textasciitilde{}1e{-}15), e a simulação dá t\_p = 4,185 s, igual ao da resolução.%
\end{itemize}%
\paragraph{Pontos de dúvida}%
\begin{itemize}%
\item%
\textit{(key)} A resolução detalhada do gabarito tem erros de digitação, que podem ser só expoentes perdidos na transcrição: '$\omega$\_e$^2$ = $\surd$($\mu$\_e g/R - $\alpha$$^2$)' (o certo é $\surd$($\mu$e$^2$g$^2$/R$^2$ - $\alpha$$^2$)), o D final com '$\surd$($\mu$\_e g/R$^2$ - $\alpha$$^2$)' (falta o quadrado em $\mu$e$\cdot$g, o que deixa a expressão dimensionalmente inconsistente), 't\_e = $\omega$\_e/a' em vez de $\omega$\_e/$\alpha$ e, na linha final de t\_P, a raiz quarta de '$\mu$\_e g/R$^2$ - $\alpha$$^2$', de novo sem o quadrado. As respostas por item do gabarito estão corretas, então isso não afeta a nota da resolução. Mesmo assim, convém não copiar essas linhas.%
\item%
\textit{(statement)} O enunciado não diz o que acontece com a velocidade horizontal quando a massa toca o solo sem quicar. A resolução e o gabarito supõem que ela se conserva e que o choque é instantâneo. A própria resolução aponta, como alternativa, um atrito impulsivo de impulso $\mu$c$\cdot$m$\cdot$$\surd$(2gh), que daria t\_p = t1 + v0/($\mu$c g). A leitura usual, adotada pelos dois, é a de conservar a velocidade.%
\item%
\textit{(solution)} Há hipóteses implícitas nos dois: (i) $\mu$e$\cdot$g > $\alpha$R (senão a massa escorrega já em t = 0, sem velocidade); (ii) como a massa está na borda, ela sai do disco no mesmo instante em que o atrito estático satura, porque a força de inércia tem componente radial para fora, e por isso o atrito cinético massa--carrossel, que o enunciado não dá, não entra. A resolução declara as duas, e são razoáveis.%
\end{itemize}%
\section*{Questão 4}%
\href{https://poliedroresolve.sistemapoliedro.com.br/vestibulares/ita/2027/ita-2-fase-dia-3-22-10-2026-2-fase-dia-3/dissertativa-4-fisica-2-fase-dia-3-ita-2027}{Resolução do gabarito}\par\smallskip%
\begin{longtable}{@{}p{0.06\textwidth} p{0.33\textwidth} p{0.33\textwidth} p{0.1\textwidth} p{0.08\textwidth}@{}}\toprule Item & Resolução & Gabarito & Veredito & Certo \\ \midrule
a & h\_f = h\_i$\cdot$((m+M)/m)\^{}\{3/5\} & h\_f = (1 + M/m)\^{}\{3/5\}$\cdot$h\_i & igual & ambos \\
b & (3/2)(m+M)h\_i{[}1 - (h\_i/h\_max)\^{}\{2/3\}{]} = m(h\_max - h\_i), equivalente a 3(m+M)h\_i{[}1 - (h\_i/h\_max)\^{}\{2/3\}{]} = 2m(h\_max - h\_i). Para M = m, h\_max $\approx$ 2,26$\cdot$h\_i & Equação geral h\_max = (3M+5m)g$\cdot$h\_i/(3p\_fA + 2mg), com p\_f indeterminado. Na interpretação 1 (adiabática reversível): 3(M+m)(h\_i/h\_max)\^{}\{5/3\}h\_max + 2m$\cdot$h\_max = (3M+5m)h\_i. Na interpretação 2 (repouso final): h\_eq = (1 + 3M/(5m))h\_i & equivalente & ambos \\
\bottomrule\end{longtable}%
\begin{itemize}%
\item%
\textbf{a)} É a mesma expressão, já que (m+M)/m = 1 + M/m. As duas resoluções usam a adiabática reversível p\_iV\_i\^{}\{5/3\} = p\_fV\_f\^{}\{5/3\}, com p\_i = (m+M)g/A e p\_f = mg/A. Confere também nos limites: com M$\to$0, h\_f$\to$h\_i, e h\_f > h\_i, como numa expansão.%
\item%
\textbf{b)} Como (h\_i/h\_max)\^{}\{5/3\}$\cdot$h\_max = h\_i$\cdot$(h\_i/h\_max)\^{}\{2/3\}, a interpretação 1 do gabarito fica 3(M+m)h\_i(h\_i/h\_max)\^{}\{2/3\} + 2m$\cdot$h\_max = 3(M+m)h\_i + 2m$\cdot$h\_i, ou seja, 3(m+M)h\_i{[}1 - (h\_i/h\_max)\^{}\{2/3\}{]} = 2m(h\_max - h\_i). É exatamente a forma equivalente da resolução; a diferença das duas equações deu zero em valores aleatórios. As duas resoluções fazem o mesmo balanço de energia: Q = 0, o pistão sai do repouso e para em h\_max, logo W\_gás = mg(h\_max - h\_i) = -$\Delta$U = (3/2)(P\_iV\_i - P\_maxV\_max). Refiz as contas. Com M = m, h\_f = 2\^{}\{0,6\}h\_i $\approx$ 1,516h\_i e a raiz não trivial é h\_max $\approx$ 2,256h\_i, como a resolução diz. A análise de f(x) também está certa: f' = kx\^{}\{-5/3\} - 1 se anula em h\_f, f'' < 0 e a raiz é única, com h\_max > h\_f. O limite M ? m dá h\_max $\approx$ 2h\_f - h\_i; para M/m = 0,01 a raiz é 1,01201 contra 1,01198 da aproximação. A interpretação 2 do gabarito (h\_eq) não responde ao item, porque calcula o repouso final depois de oscilações amortecidas, não a altura máxima. Para essa outra pergunta, porém, a conta está certa: 3mg$\cdot$h + 2mg$\cdot$h = (3M+5m)g$\cdot$h\_i.%
\end{itemize}%
\paragraph{Pontos de dúvida}%
\begin{itemize}%
\item%
\textit{(statement)} No item b o enunciado não diz se o gás evolui de forma reversível depois da retirada instantânea da areia; o próprio Poliedro aponta essa falta. A resolução declara a hipótese de que o gás é uniforme e isentrópico (pistão lento em comparação com a velocidade do som; razão entre os tempos \textasciitilde{} $\surd$(m\_gás/m) ? 1). Essa hipótese coincide com a interpretação 1 do gabarito e é a mais coerente com o pedido de 'equação que relaciona h\_max'. Ainda assim, convém conferir o critério oficial do ITA: se ele adotar outra leitura, a nota do item b pode mudar.%
\item%
\textit{(key)} A transcrição do gabarito perdeu a formatação. Aparece '(M + mg)/A' onde deveria ser '(M+m)g/A', e '1 + M/m\^{}(3/5)' onde deveria ser '(1+M/m)\^{}\{3/5\}'. No item b, h\_f é usado no lugar de h\_max. Os resultados finais não mudam.%
\item%
\textit{(key)} A interpretação 2 do gabarito, h\_eq = (1 + 3M/(5m))h\_i, dá a altura de equilíbrio final depois de oscilações dissipadas. Isso não é a altura máxima pedida. A conta está certa para essa outra pergunta, mas o resultado não serve como resposta do item b.%
\end{itemize}%
\section*{Questão 5}%
\href{https://poliedroresolve.sistemapoliedro.com.br/vestibulares/ita/2027/ita-2-fase-dia-3-22-10-2026-2-fase-dia-3/dissertativa-5-fisica-2-fase-dia-3-ita-2027}{Resolução do gabarito}\par\smallskip%
\begin{longtable}{@{}p{0.06\textwidth} p{0.33\textwidth} p{0.33\textwidth} p{0.1\textwidth} p{0.08\textwidth}@{}}\toprule Item & Resolução & Gabarito & Veredito & Certo \\ \midrule
a & p' = -n2pR/{[}n1R + (n2-n1)p{]}: imagem virtual no meio n1, a |p'| = n2pR/{[}n1R + (n2-n1)p{]} à esquerda do vértice. Alturas y' = $\pm$ n1Rh/{[}n1R + (n2-n1)p{]}, imagens direitas e reduzidas. & p' = -n2pR/(n1R + (n2-n1)p), imagens virtuais no meio n1; y' = $\pm$ n1Rh/(n1R + (n2-n1)p) & igual & ambos \\
b & $\Delta$y = $\lambda$(D + n2p/n1)/(2n2h) $\approx$ $\lambda$D/(2n2h), com $\lambda$ no vácuo. Se $\lambda$ for o comprimento de onda no meio n1, fica n1$\lambda$D/(2n2h). & y = $\lambda$D/(2n2h), com $\lambda$ no vácuo & equivalente & ambos \\
\bottomrule\end{longtable}%
\begin{itemize}%
\item%
\textbf{a)} Refiz pela equação do dioptro com R\_alg = -R, já que o centro C fica do lado da luz incidente: p' = -n2pR/{[}n1R+(n2-n1)p{]}. O aumento m = -n1p'/(n2p) dá n1R/{[}n1R+(n2-n1)p{]}. A resolução chega ao mesmo valor por outro caminho, usando o raio que passa por C (h(R-|p'|)/(R-p)). Conferi numericamente com n1=1, n2=1,5, R=0,3 e p=0,1: p' = -0,12857 e m = 0,857 pelos dois caminhos. O sinal, a natureza virtual da imagem e as alturas coincidem com o gabarito.%
\item%
\textbf{b)} A resolução mantém o termo de correção n2p/n1, que vem das fontes virtuais a n2p/n1 do vértice. Como o enunciado impõe p ? D, a resposta se reduz a $\lambda$D/(2n2h), que é a do gabarito, e a própria resolução escreve essa forma aproximada. Fiz um traçado exato de caminho óptico (Fermat), com $\lambda$=600 nm, D=1 m, h=0,5 mm, p=1 cm, n1=1 e n2=1,5. A primeira franja ficou em y = 0,4060 mm, contra 0,4060 mm pela fórmula com correção e 0,4000 mm pela aproximada. Isso confirma as duas formas dentro da hipótese p ? D. As duas adotam $\lambda$ como comprimento de onda no vácuo.%
\end{itemize}%
\paragraph{Pontos de dúvida}%
\begin{itemize}%
\item%
\textit{(statement)} O enunciado diz apenas que as fontes, imersas no meio n1, têm 'comprimento de onda $\lambda$', sem dizer se é no vácuo ou no meio n1. O gabarito e a resolução tomam $\lambda$ no vácuo e chegam a $\lambda$D/(2n2h). Se $\lambda$ for o comprimento de onda no meio n1, a resposta vira n1$\lambda$D/(2n2h). A resolução cita essa alternativa, mas o gabarito não.%
\item%
\textit{(solution)} No item b, a resposta principal da resolução é $\lambda$(D + n2p/n1)/(2n2h), mais precisa que a do gabarito, e a forma $\approx$$\lambda$D/(2n2h) vem em seguida. Como o enunciado dá p ? D, o esperado é $\lambda$D/(2n2h). Não deve haver desconto, porque a resolução declara a aproximação.%
\item%
\textit{(key)} No item b, o gabarito usa sen $\theta$2 $\approx$ y/D e ignora o deslocamento p$\cdot$tan $\theta$1 no meio n1. Isso vale pela hipótese p ? D e não é erro, só uma aproximação que o gabarito não declara.%
\item%
\textit{(statement)} O item a pede as posições verticais 'em função de n1, n2, R e p', mas elas dependem também de h, como nas duas respostas. É uma omissão de redação do enunciado.%
\end{itemize}%
\section*{Questão 6}%
\href{https://poliedroresolve.sistemapoliedro.com.br/vestibulares/ita/2027/ita-2-fase-dia-3-22-10-2026-2-fase-dia-3/dissertativa-6-fisica-2-fase-dia-3-ita-2027}{Resolução do gabarito}\par\smallskip%
\begin{longtable}{@{}p{0.06\textwidth} p{0.33\textwidth} p{0.33\textwidth} p{0.1\textwidth} p{0.08\textwidth}@{}}\toprule Item & Resolução & Gabarito & Veredito & Certo \\ \midrule
a & q = 3C$\varepsilon$/4 (C' = 3C, V = $\varepsilon$/4) & q = 3C$\varepsilon$/4 (C' = 3C, V3 = $\varepsilon$/4) & igual & ambos \\
b & compressão $\Delta$x = 2d/3 (separação final d/3) & compressão $\Delta$x = 2d/3 (distância final d/3) & igual & ambos \\
c & k = 27C$\varepsilon$$^2$/(64d$^2$) = 27$\varepsilon$$_{0}$a$^2$$\varepsilon$$^2$/(64d$^3$) & k = 27C$\varepsilon$$^2$/(64d$^2$) & igual & ambos \\
\bottomrule\end{longtable}%
\begin{itemize}%
\item%
\textbf{a)} Conservação de carga no conjunto isolado capacitor + placas, com os dois em paralelo: C$\varepsilon$ = (C + C')$\cdot$$\varepsilon$/4, logo C' = 3C e q = C'$\cdot$$\varepsilon$/4 = 3C$\varepsilon$/4. Confirmei com frações exatas.%
\item%
\textbf{b)} Com C(x) = $\varepsilon$$_{0}$A/x = Cd/x e C' = 3C, a separação final é x = d/3. A mola começa no comprimento natural, então fica comprimida de 2d/3. O sentido (compressão) também coincide.%
\item%
\textbf{c)} A força entre as placas é F = q$^2$/(2$\varepsilon$$_{0}$A). Com $\varepsilon$$_{0}$A = Cd, F = (9C$^2$$\varepsilon$$^2$/16)/(2Cd) = 9C$\varepsilon$$^2$/(32d). Do equilíbrio k$\cdot$(2d/3) = F sai k = 27C$\varepsilon$$^2$/(64d$^2$). A forma em $\varepsilon$$_{0}$a$^2$/d$^3$ é só a mesma expressão reescrita. Conferido com frações exatas: F = 9/32 e k = 27/64 em unidades C = $\varepsilon$ = d = 1.%
\end{itemize}%
\paragraph{Pontos de dúvida}%
\begin{itemize}%
\item%
\textit{(statement)} Com o k obtido, o equilíbrio em t$_{3}$ fica exatamente no limite de estabilidade. Como a carga total C$\varepsilon$ se conserva, a força elétrica em função da separação x é F(x) = C$\varepsilon$$^2$d/{[}2(x+d)$^2${]}, e |dF/dx| em x = d/3 vale 27C$\varepsilon$$^2$/(64d$^2$), que é igual a k. A força resultante tem raiz dupla em x = d/3 e puxa as placas para dentro dos dois lados desse ponto. Na prática o equilíbrio é semiestável: com dinâmica sem amortecimento as placas passariam do ponto e colapsariam. A resolução comenta isso; o gabarito não. Não muda as respostas, mas é uma fragilidade física do enunciado que vale mostrar ao professor.%
\item%
\textit{(solution)} Hipótese implícita nas duas resoluções: as placas estão penduradas por fios, mas a componente horizontal da tração (efeito de pêndulo) é desprezada. O enunciado não dá massa nem comprimento dos fios, então isso é inevitável. A resolução declara a hipótese; o gabarito a usa sem dizer.%
\item%
\textit{(key)} O texto do gabarito tem frações e linhas grudadas, provavelmente por perda de formatação na extração (ex.: 'C$\varepsilon$ = C + C'$\varepsilon$/4C' = 3C', 'C' = $\varepsilon$\_0A/x = 3$\varepsilon$\_0A/dx = d/3'). Os resultados finais estão corretos e coerentes; é só redação ou transcrição.%
\end{itemize}%
\section*{Questão 7}%
\href{https://poliedroresolve.sistemapoliedro.com.br/vestibulares/ita/2027/ita-2-fase-dia-3-22-10-2026-2-fase-dia-3/dissertativa-7-fisica-2-fase-dia-3-ita-2027}{Resolução do gabarito}\par\smallskip%
\begin{longtable}{@{}p{0.06\textwidth} p{0.33\textwidth} p{0.33\textwidth} p{0.1\textwidth} p{0.08\textwidth}@{}}\toprule Item & Resolução & Gabarito & Veredito & Certo \\ \midrule
a & |V\_o| = (V\_s/4)$\cdot$|G$\varepsilon$ + $\alpha$$\Delta$T|. O sinal é + com o extensômetro em R2 ou R3 e - em R1 ou R4. & |V0| = (Vs/4)$\cdot$|G$\varepsilon$ + $\alpha$$\Delta$T| & igual & ambos \\
b & Os quatro dispositivos devem ser extensômetros: R2 e R3 em pontos de tração e R1 e R4 em pontos de compressão (ou o inverso). Assim, |V\_o| = V\_s$\cdot$G$\cdot$|$\varepsilon$|, sem efeito térmico e com 4 vezes a sensibilidade do item a. & Quatro extensômetros: R2 e R3 na tração, R1 e R4 na compressão (ou tudo trocado), dando V0 = Vs$\cdot$G$\cdot$$\varepsilon$, sem efeito térmico. & igual & ambos \\
\bottomrule\end{longtable}%
\begin{itemize}%
\item%
\textbf{a)} Refiz a conta. Com R\_i = R0(1+x\_i), a saída é V\_o = V\_D - V\_B $\approx$ (V\_s/4)(x2 + x3 - x1 - x4). Com um só extensômetro, só um x\_k vale G$\varepsilon$+$\alpha$$\Delta$T, então |V\_o| = V\_s|G$\varepsilon$+$\alpha$$\Delta$T|/4. Conferi também numericamente, com V\_s=10 V, G=2,1, $\alpha$=1,2e{-}5, $\Delta$T=7 e $\varepsilon$=3e{-}4: o valor exato é 1,7844 mV e a aproximação linear dá 1,7850 mV, com o sinal mudando conforme o braço como a resolução diz. A resolução e o gabarito chegam à mesma expressão e à mesma dedução.%
\item%
\textbf{b)} A resolução e o gabarito propõem a mesma configuração e chegam à mesma saída. A conta exata, com x2=x3=G$\varepsilon$+$\alpha$$\Delta$T e x1=x4=-G$\varepsilon$+$\alpha$$\Delta$T, dá V\_o = V\_sG$\varepsilon$/(1+$\alpha$$\Delta$T) $\approx$ V\_sG$\varepsilon$. O termo térmico se cancela e só se despreza o produto $\varepsilon$$\alpha$$\Delta$T, que o enunciado manda desprezar. No teste numérico, o valor exato foi 6,29947 mV contra V\_sG$\varepsilon$ = 6,3 mV. As duas aceitam a troca de todos os pontos (V\_o = -V\_sG$\varepsilon$). A resolução escreve |V\_o| = V\_sG|$\varepsilon$| em módulo e o gabarito escreve V0 = VsG$\varepsilon$; a diferença é só de notação.%
\end{itemize}%
\paragraph{Pontos de dúvida}%
\begin{itemize}%
\item%
\textit{(statement)} O item b pede para indicar, para cada dispositivo, se ele deve ser extensômetro ou resistor R0. Isso pode levar o aluno a responder com meia ponte: dois extensômetros adjacentes, um na tração e outro na compressão, mais dois resistores. Essa montagem também cancela o efeito térmico, mas tem metade da sensibilidade (V\_sG$\varepsilon$/2). Como o enunciado pede 'máxima precisão', a ponte completa é a resposta ótima, como dizem o gabarito e a resolução. Vale confirmar com o professor quanto crédito a meia ponte receberia.%
\item%
\textit{(statement)} No item b, $\varepsilon$ é tratado como a intensidade comum da deformação: +$\varepsilon$ nos pontos de tração e -$\varepsilon$ nos de compressão. O enunciado define $\varepsilon$>0 para tração. A resolução escreve |V\_o| = V\_sG|$\varepsilon$| e o gabarito escreve V0 = VsG$\varepsilon$. A diferença é só de convenção e não muda a resposta.%
\item%
\textit{(key)} Na transcrição da resolução do gabarito faltam colchetes e há uma passagem com '1/2(1 + x\_3)1 + (x\_1 + x\_3)/2\^{}({-}1)'. São defeitos de extração do texto, não erros físicos. O resultado final (1/2 + (x3-x1)/4) está correto.%
\end{itemize}%
\section*{Questão 8}%
\href{https://poliedroresolve.sistemapoliedro.com.br/vestibulares/ita/2027/ita-2-fase-dia-3-22-10-2026-2-fase-dia-3/dissertativa-8-fisica-2-fase-dia-3-ita-2027}{Resolução do gabarito}\par\smallskip%
\begin{longtable}{@{}p{0.06\textwidth} p{0.33\textwidth} p{0.33\textwidth} p{0.1\textwidth} p{0.08\textwidth}@{}}\toprule Item & Resolução & Gabarito & Veredito & Certo \\ \midrule
a & v = $\surd$(2qV/m) & v = $\surd$(2qV/m) & igual & ambos \\
b & Movimento helicoidal uniforme: MRU ao longo de B com v$\cdot$cos$\theta$, composto com MCU no plano perpendicular a B, de raio R = m$\cdot$v$\cdot$sen$\theta$/(qB) e período T = 2$\pi$m/(qB). A justificativa é que F = q v$\times$B é perpendicular a B (logo v? é constante) e perpendicular a v (logo |v| é constante e F atua como resultante centrípeta). & Movimento helicoidal uniforme: MCU no plano perpendicular a B (R = m v sen$\theta$/(qB), T = 2$\pi$m/(qB)) composto com MU ao longo de B (v cos$\theta$) & igual & ambos \\
c & cos$\theta$ = qBL/(2$\pi$n$\cdot$m$\cdot$v) = (BL/(2$\pi$n))$\cdot$$\surd$(q/(2mV)), n = 1, 2, 3, \ldots{} (com cos$\theta$ < 1) & cos$\theta$ = qBL/(2$\pi$ n m v) = (BL/(2$\pi$ n))$\cdot$$\surd$(q/(2mV)), n $\in$ $\mathbb{N}$* (com cos$\theta$ $\leq$ 1) & igual & ambos \\
\bottomrule\end{longtable}%
\begin{itemize}%
\item%
\textbf{a)} Os dois usam o teorema da energia cinética, qV = mv$^2$/2, partindo do repouso. O resultado é idêntico.%
\item%
\textbf{b)} A classificação, a justificativa, R e T são os mesmos. A resolução também dá o passo p = 2$\pi$m v cos$\theta$/(qB) e cita os casos{-}limite $\theta$ = 0 (MRU) e $\theta$ = 90$^\circ$ (MCU). Nada disso contradiz o gabarito.%
\item%
\textbf{c)} A expressão é a mesma e n também é inteiro positivo nos dois. A única diferença está na restrição: a resolução escreve cos$\theta$ < 1 e o gabarito, cos$\theta$ $\leq$ 1. Para $\theta$ = 0 a partícula nem sai da linha de campo, então não há 'retorno'. Por isso a desigualdade estrita da resolução é até mais coerente com o enunciado, e essa diferença não muda a resposta. Conferi numericamente o exemplo da resolução (próton, B = 0,5 T, V = 1000 V, $\theta$ = 0,7 rad, n = 3): L $\approx$ 0,1317 m, e a fórmula devolve cos$\theta$ = 0,7648 = cos(0,7).%
\end{itemize}%
\paragraph{Pontos de dúvida}%
\begin{itemize}%
\item%
\textit{(solution)} Na restrição sobre n, a resolução usa cos$\theta$ < 1 (exclui $\theta$ = 0) e o gabarito usa cos$\theta$ $\leq$ 1. Para $\theta$ = 0 o movimento é um MRU sobre a própria linha, sem retornos de fato, e a figura mostra $\theta$ $\neq$ 0. A diferença é só de redação e não afeta a nota.%
\item%
\textit{(key)} A observação final do gabarito diz que 'com n = 1 a partícula chega a 2 no primeiro retorno'. Isso só é possível quando BL/(2$\pi$)$\cdot$$\surd$(q/(2mV)) $\leq$ 1, o que não fica explícito. Isso é imprecisão de redação, não erro na resposta.%
\end{itemize}%
\section*{Questão 9}%
\href{https://poliedroresolve.sistemapoliedro.com.br/vestibulares/ita/2027/ita-2-fase-dia-3-22-10-2026-2-fase-dia-3/dissertativa-9-fisica-2-fase-dia-3-ita-2027}{Resolução do gabarito}\par\smallskip%
\begin{longtable}{@{}p{0.06\textwidth} p{0.33\textwidth} p{0.33\textwidth} p{0.1\textwidth} p{0.08\textwidth}@{}}\toprule Item & Resolução & Gabarito & Veredito & Certo \\ \midrule
a & E = E0{[}1+({-}1)\^{}n{]} cos(kx - $\omega$t) $\hat z$ = 2E0 cos$^2$(n$\pi$/2) cos(kx - $\omega$t) $\hat z$. A amplitude é E0{[}1+({-}1)\^{}n{]} e a onda se propaga com velocidade $\omega$/k = c no sentido +x. & E = E0{[}1 + (-1)\^{}n{]} cos(kx - $\omega$t) $\hat z$; amplitude E0{[}1 + (-1)\^{}n{]}; onda no sentido +x & igual & ambos \\
b & n par: interferência construtiva, amplitude 2E0 e intensidade 4 vezes a de uma onda isolada. n ímpar: interferência destrutiva total, E = 0. & n par: construtiva (amplitude 2E0); n ímpar: destrutiva total (amplitude 0) & igual & ambos \\
c & B = -B0 cos(kx - $\omega$t) $\hat y$, com B0 = E0{[}1+(-1)\^{}n{]}/c. B oscila ao longo de y, em fase com E, e E$\times$B aponta para +x. & B = -B0 cos(kx - $\omega$t) $\hat y$ (E $\times$ B no sentido +x?) & igual & ambos \\
\bottomrule\end{longtable}%
\begin{itemize}%
\item%
\textbf{a)} Mesma expressão. A resolução chega a ela pela soma em produto, 2E0 cos(n$\pi$/2) cos($\varphi$ + n$\pi$/2), e confere direto com cos($\varphi$ + n$\pi$) = (-1)\^{}n cos $\varphi$, que é o caminho do gabarito. A forma 2E0 cos$^2$(n$\pi$/2) é idêntica a E0{[}1+(-1)\^{}n{]}. O sentido +x vem da fase kx - $\omega$t ser constante em x = x0 + ($\omega$/k)t, o que está correto. O enunciado foi conferido no PDF original (página 10 impressa) e bate com a transcrição.%
\item%
\textbf{b)} Conclusão idêntica. A resolução acrescenta, de forma correta, que d = n$\lambda$/2 e que a intensidade é 4 vezes a de uma onda isolada no caso construtivo.%
\item%
\textbf{c)} Mesmo vetor e mesma orientação. A resolução chega ao resultado pela lei de Faraday: (?$\times$E)\_y = -$\partial$E\_z/$\partial$x, logo $\partial$B\_y/$\partial$t = $\partial$E\_z/$\partial$x = -kE\_n sen(kx-$\omega$t), e daí B\_y = -(k/$\omega$)E\_n cos(kx-$\omega$t). O sinal está correto. O gabarito chega ao mesmo vetor exigindo que E$\times$B aponte para +x? ($\hat z$ $\times$ (-$\hat y$) = x?). A resolução também dá o valor B0 = E0{[}1+(-1)\^{}n{]}/c, coerente com o item a, e oferece uma leitura alternativa (B0 como amplitude de cada onda isolada), mas a resposta principal coincide com a do gabarito.%
\end{itemize}%
\paragraph{Pontos de dúvida}%
\begin{itemize}%
\item%
\textit{(solution)} No item c a resolução apresenta duas leituras de B0: a amplitude resultante, que é a principal e coincide com o gabarito, e a amplitude de cada onda isolada. A resposta principal está certa, mas a alternativa pode ser vista como hesitação. O enunciado diz que o campo resultante 'possui amplitude de módulo B0', o que favorece a primeira leitura.%
\item%
\textit{(statement)} Para n ímpar o campo resultante é identicamente nulo. Dizer que ele 'se propaga no sentido +x' (item a) e ter B = -B0 cos(...) $\hat y$ com B0 = 0 (item c) só vale de forma trivial. Nem a resolução nem o gabarito tratam isso como problema, e não muda a nota.%
\end{itemize}%
\section*{Questão 10}%
\href{https://poliedroresolve.sistemapoliedro.com.br/vestibulares/ita/2027/ita-2-fase-dia-3-22-10-2026-2-fase-dia-3/dissertativa-10-fisica-2-fase-dia-3-ita-2027}{Resolução do gabarito}\par\smallskip%
\begin{longtable}{@{}p{0.06\textwidth} p{0.33\textwidth} p{0.33\textwidth} p{0.1\textwidth} p{0.08\textwidth}@{}}\toprule Item & Resolução & Gabarito & Veredito & Certo \\ \midrule
a & $\nu$' = $\nu$ / {[}1 + (h$\nu$/(m\_e c$^2$))(1 - cos $\theta$){]} & $\nu$' = $\nu$/(1 + (h$\nu$/(m\_e c$^2$))(1 - cos$\theta$)) (Compton) & igual & ambos \\
b & |V| = c($\nu$$^2$ - $\nu$'$^2$)/($\nu$$^2$ + $\nu$'$^2$) = c$\cdot$x(2+x)/(x$^2$+2x+2), com x = (h$\nu$/(m\_e c$^2$))(1 - cos $\theta$); O se move na direção do fóton espalhado, em sentido oposto ao de propagação dele (ao seu encontro) & V = c$\cdot${[}(1 + h$\nu$(1-cos$\theta$)/(m\_e c$^2$))$^2$ - 1{]}/{[}(1 + h$\nu$(1-cos$\theta$)/(m\_e c$^2$))$^2$ + 1{]}, com $\beta$ = ($\nu$$^2$ - $\nu$'$^2$)/($\nu$$^2$ + $\nu$'$^2$); observador na direção do fóton espalhado, em sentido oposto ao de propagação (aproximando{-}se) & equivalente & ambos \\
\bottomrule\end{longtable}%
\begin{itemize}%
\item%
\textbf{a)} As duas expressões são idênticas: é a fórmula de Compton, obtida pela conservação da energia relativística e do momento, com a lei dos cossenos no triângulo dos momentos. A dedução da resolução está correta, igual à do gabarito: os termos h$^2$$\nu$$^2$, h$^2$$\nu$'$^2$ e m\_e$^2$c$^{4}$ se cancelam e sobra m\_ec$^2$($\nu$-$\nu$') = h$\nu$$\nu$'(1-cos$\theta$).%
\item%
\textbf{b)} É a mesma resposta escrita de outro jeito: (1+x)$^2$ - 1 = x(2+x) e (1+x)$^2$ + 1 = x$^2$ + 2x + 2. A segunda forma da resolução, {[}m\_ec$^2$ + h$\nu$(1-cos$\theta$){]}$^2$ - (m\_ec$^2$)$^2$ sobre {[}m\_ec$^2$ + h$\nu$(1-cos$\theta$){]}$^2$ + (m\_ec$^2$)$^2$, vira a do gabarito quando se divide numerador e denominador por (m\_ec$^2$)$^2$. Conferi numericamente: as expressões coincidem e, com $\nu$\_O = $\nu$'$\surd$((1+$\beta$)/(1-$\beta$)), resulta $\nu$\_O = $\nu$. O sentido também coincide: os dois põem o observador indo ao encontro do fóton espalhado, o que dá o desvio para o azul. A resolução usa $\beta$ algébrico ($\beta$ < 0 no sentido de propagação) e apresenta |V| junto com o sentido. O gabarito usa $\beta$ > 0 como rapidez de aproximação. As duas convenções são consistentes. Também conferi a estimativa da resolução (h$\nu$ = 17,5 keV, $\theta$ = 90$^\circ$): x $\approx$ 0,0342 e |V| $\approx$ 0,0337c $\approx$ 1,0$\times$10$^{7}$ m/s.%
\end{itemize}%
\paragraph{Pontos de dúvida}%
\begin{itemize}%
\item%
\textit{(key)} A resolução do gabarito tem erros de digitação que não afetam a resposta final. Na lei dos cossenos aparece "p\_e$^2$ = h$\nu$/c$^2$ + (h$\nu$')/c$^2$ - \ldots{}", quando deveria ser (h$\nu$/c)$^2$ + (h$\nu$'/c)$^2$ - \ldots{}. Nas últimas linhas do item b faltam parênteses: "1 + h$\nu$/(m\_ec$^2$)(1-cos$\theta$)$^2$" deveria ser {[}1 + h$\nu$(1-cos$\theta$)/(m\_ec$^2$){]}$^2$. A resposta{-}resumo do gabarito traz os colchetes corretamente.%
\item%
\textit{(statement)} O enunciado diz apenas que O se move "na mesma direção do fóton espalhado", sem dizer o sentido. A resposta completa precisa indicar que O vai ao encontro do fóton (sentido oposto ao de propagação), porque $\nu$' < $\nu$ exige desvio para o azul. A resolução e o gabarito dizem isso explicitamente. Uma resposta só com o módulo de V poderia perder parte da pontuação.%
\item%
\textit{(solution)} A resolução chama a colisão de "elástica" no sentido de que as massas de repouso se conservam. A energia cinética total também se conserva (o fóton não tem massa e o elétron mantém m\_e), então a interpretação está certa. É uma hipótese implícita razoável, que o gabarito também adota.%
\end{itemize}%
\end{document}